\documentclass[UTF8,fontset=none,12pt]{ctexart}

\setCJKmainfont{Songti SC}
\setCJKsansfont{PingFang SC}
\setCJKmonofont{Songti SC}

\usepackage[a4paper,margin=2.35cm]{geometry}
\usepackage{amsmath,amssymb,amsthm,mathtools,bm}
\usepackage{graphicx}
\usepackage{float}
\usepackage{booktabs}
\usepackage{enumitem}
\usepackage[hidelinks]{hyperref}

\newcommand{\C}{\mathbb{C}}
\newcommand{\epsmach}{\varepsilon_{\mathrm{machine}}}
\newcommand{\fl}{\operatorname{fl}}
\theoremstyle{definition}
\newtheorem{theorem}{定理}
\renewcommand{\thetheorem}{15.\arabic{theorem}}
\renewcommand{\thefigure}{15.\arabic{figure}}

\title{数值线性代数\quad 第 15 讲：稳定性漫谈}
\author{Lloyd N. Trefethen、David Bau III\quad 中文阅读版}
\date{}

\begin{document}
\maketitle
\pagestyle{plain}

我们继续讨论稳定性，考察稳定与不稳定算法的例子.随后将讨论一个把良态性和稳定性联系起来的基本思想：后向误差分析.自 20 世纪 50 年代以来，它的力量已经在无数应用中得到证明.

\section*{浮点运算的稳定性}

最简单的四个计算问题是 $+$、$-$、$\times$ 和 $\div$.在算法选择方面没有太多可说的!当然，我们通常会使用计算机提供的浮点运算 $\oplus$、$\ominus$、$\otimes$ 和 $\oslash$.事实证明，公理 (13.5) 和 (13.7) 蕴含：这四个典型算法都是后向稳定的.

下面以减法为例说明，因为这是人们最容易担心不稳定的基本运算.与例 12.3 一样，数据空间 $X$ 是二维向量空间 $\C^2$，解空间 $Y$ 是标量空间 $\C$.由定理 14.1，我们无需指定这些空间上的范数.

对数据 $x=(x_1,x_2)^*\in X$，减法问题对应于函数
\[
    f(x_1,x_2)=x_1-x_2,
\]
所考虑的算法可以写成
\begin{equation*}
    \widetilde f(x_1,x_2)=\fl(x_1)\ominus\fl(x_2).
\end{equation*}

这个等式表示：先把 $x_1$ 和 $x_2$ 舍入为浮点数，再应用浮点减法.由式 (13.5)，有
\[
    \fl(x_1)=x_1(1+\epsilon_1),
    \qquad
    \fl(x_2)=x_2(1+\epsilon_2),
    \qquad
    |\epsilon_1|,|\epsilon_2|\le\epsmach.
\]
由式 (13.7)，有
\[
    \fl(x_1)\ominus\fl(x_2)
    =\bigl(\fl(x_1)-\fl(x_2)\bigr)(1+\epsilon_3),
    \qquad
    |\epsilon_3|\le\epsmach.
\]
合并这些等式得到
\begin{align*}
    \fl(x_1)\ominus\fl(x_2)
    &=[x_1(1+\epsilon_1)-x_2(1+\epsilon_2)](1+\epsilon_3)\\
    &=x_1(1+\epsilon_1)(1+\epsilon_3)
      -x_2(1+\epsilon_2)(1+\epsilon_3)\\
    &=x_1(1+\epsilon_4)-x_2(1+\epsilon_5),
\end{align*}
其中
\[
    |\epsilon_4|,|\epsilon_5|
    \le 2\epsmach+O(\epsmach^2)
\]
（见习题 14.2）.换句话说，计算结果
\[
    \widetilde f(x)=\fl(x_1)\ominus\fl(x_2)
\]
恰好等于 $\widetilde x_1-\widetilde x_2$，其中 $\widetilde x_1,\widetilde x_2$ 满足
\[
    \frac{|\widetilde x_1-x_1|}{|x_1|}
    =O(\epsmach),
    \qquad
    \frac{|\widetilde x_2-x_2|}{|x_2|}
    =O(\epsmach).
\]
隐含常数取任意 $C>2$ 都足够.对 $\C^2$ 上的任意范数，这就推出式 (14.5).

\section*{更多例子}

\subsection*{例 15.1：内积}

给定向量 $x,y\in\C^m$，希望计算内积
\[
    \alpha=x^*y.
\]
显然的算法是：使用 $\otimes$ 计算逐项乘积 $\widetilde x_i\otimes\widetilde y_i$，再使用 $\oplus$ 将它们相加，得到计算结果 $\widetilde\alpha$.可以证明，这个算法是后向稳定的；第 17 讲会隐含地完成这一证明.\hfill$\square$

\subsection*{例 15.2：外积}

另一方面，设希望对向量 $x\in\C^m$、$y\in\C^n$ 计算秩一外积
\[
    A=xy^*.
\]
显然的算法是：使用 $\otimes$ 计算 $mn$ 个乘积 $\widetilde x_i\otimes\widetilde y_j$，并将它们收集为矩阵 $\widetilde A$.这个算法是稳定的，但不是后向稳定的.原因在于，矩阵 $\widetilde A$ 极不可能恰好具有秩 1，因此一般不能写成
\[
    \widetilde A=(x+\delta x)(y+\delta y)^*.
\]
一般而言，对于解空间 $Y$ 的维数大于问题空间 $X$ 维数的问题，后向稳定性是罕见的.\hfill$\square$

\subsection*{例 15.3}

假设使用 $\oplus$ 计算 $x+1$，其中 $x\in\C$：
\[
    \widetilde f(x)=\fl(x)\oplus1.
\]
这个算法稳定，但不是后向稳定的.原因是，当 $x\approx0$ 时，加法 $\oplus$ 会引入大小为 $O(\epsmach)$ 的绝对误差.相对于 $x$ 的大小，这些误差是无界的，因此不能解释为数据发生了小的相对扰动.这个例子说明，后向稳定性是一种相当特殊的性质：在某些情形下它是合理目标，在另一些情形下则不是.注意，如果问题是对数据 $x,y$ 计算 $x+y$，那么对应算法就是后向稳定的.\hfill$\square$

\subsection*{例 15.4}

对于计算 $\sin x$ 或 $\cos x$ 的计算机程序或计算器，合理的期待是什么？答案仍然是稳定性，而不是后向稳定性.对于 $\cos x$，这来自前一个例子中的事实 $\cos0\ne0$.对于 $\sin x$ 和 $\cos x$ 二者，函数在某些点导数为零这一事实也排除了后向稳定性.

例如，设计算机计算
\[
    f(x)=\sin x,\qquad x=\frac{\pi}{2}-\delta,\qquad 0<\delta\ll1.
\]
假设我们足够幸运，得到的计算结果正好是正确答案在浮点系统中的舍入值：
\[
    \widetilde f(x)=\fl(\sin x).
\]
由于
\[
    f'(x)=\cos x\approx\delta,
\]
所以存在某个 $\widetilde x$ 使 $\widetilde f(x)=f(\widetilde x)$，并且
\[
    \widetilde x-x
    \approx\frac{\widetilde f(x)-f(x)}{\delta}
    =O\left(\frac{\epsmach}{\delta}\right).
\]
由于 $\delta$ 可以任意小，这个后向误差不可能具有 $O(\epsmach)$ 的量级.\hfill$\square$

\section*{一个不稳定算法}

以上都是玩具例子.下面给出一个更实质的例子：使用特征多项式求矩阵的特征值.

由于 $z$ 是 $A$ 的特征值当且仅当 $p(z)=0$，其中
\[
    p(z)=\det(zI-A)
\]
是特征多项式，所以 $p$ 的根就是 $A$ 的特征值（第 24 讲）.这启发出如下的特征值计算方法：
\begin{enumerate}
    \item 求出特征多项式的系数；
    \item 求这个多项式的根.
\end{enumerate}

这个算法不仅后向不稳定，而且是不稳定的，因此不应使用.即使求特征值本身是良态问题，它也可能产生远大于 $\epsmach$ 的相对误差.

不稳定性出现在第二步的求根过程中.正如例 12.5 所见，给定多项式系数后求根通常是病态问题.因此，特征多项式系数中的小误差在求根时往往会被放大，即使求根过程本身可以做到完全精确.

例如，设 $A=I$ 是 $2\times2$ 单位矩阵. $A$ 的特征值对矩阵元素的扰动不敏感，稳定算法应当能够以 $O(\epsmach)$ 的误差计算它们.然而，上述算法产生的误差量级是 $\sqrt{\epsmach}$.原因是特征多项式
\[
    p(x)=x^2-2x+1,
\]
正如例 12.5 所示，其系数计算预计会产生 $O(\epsmach)$ 量级的误差，而这些误差会使根发生 $\sqrt{\epsmach}$ 量级的变化.例如，如果 $\epsmach=10^{-16}$，计算得到的特征多项式的根可能偏离真实特征值约 $10^{-8}$，也就是损失 8 位准确数字.

在亲自尝试这个计算之前，还需要更诚实地说明一点.如果直接用上述算法计算 $2\times2$ 单位矩阵的特征值，你很可能发现完全没有误差，因为 $x^2-2x+1$ 的系数和根都是计算机能够精确表示的小整数.但是，如果对一个略有扰动的矩阵进行实验，例如
\[
    A=
    \begin{bmatrix}
        1+10^{-14}&0\\
        0&1
    \end{bmatrix},
\]
计算得到的特征值就会与真实特征值相差预期的 $\sqrt{\epsmach}$ 量级.试试看!

\section*{后向稳定算法的准确性}

设我们有一个用于问题 $f:X\to Y$ 的后向稳定算法 $\widetilde f$.它给出的结果会准确吗？答案取决于问题的条件数 $\kappa=\kappa(x)$.如果 $\kappa(x)$ 较小，结果在相对意义下会很准确；但如果 $\kappa(x)$ 很大，准确性就会按比例下降.

\begin{theorem}
设后向稳定算法用于求解问题 $f:X\to Y$，问题在数据 $x$ 处的条件数为 $\kappa=\kappa(x)$.假设计算机满足公理 (13.5) 和 (13.7)，则相对误差满足
\begin{equation}
    \frac{\|\widetilde f(x)-f(x)\|}{\|f(x)\|}
    =O\bigl(\kappa(x)\epsmach\bigr).
    \tag{15.1}
\end{equation}
\end{theorem}

\textbf{证明.}由后向稳定性的定义 (14.5)，存在某个 $\widetilde x\in X$，使得
\[
    \widetilde f(x)=f(\widetilde x),
    \qquad
    \frac{\|\widetilde x-x\|}{\|x\|}
    =O(\epsmach).
\]
由 $\kappa(x)$ 的定义 (12.5)，这意味着
\begin{equation}
    \frac{\|\widetilde f(x)-f(x)\|}{\|f(x)\|}
    \le
    \bigl(\kappa(x)+o(1)\bigr)
    \frac{\|\widetilde x-x\|}{\|x\|},
    \tag{15.2}
\end{equation}
其中 $o(1)$ 表示当 $\epsmach\to0$ 时趋于零的量.合并这两个界就得到式 (15.1). \hfill$\square$

\section*{后向误差分析}

刚才证明定理 15.1 的过程称为后向误差分析.我们通过两个步骤得到准确度估计.一步是研究问题的条件性，另一步是研究算法的稳定性.结论是：如果算法稳定，那么最终准确度就会反映问题的条件数.

从数学上说，这个过程很直接，但如果让一个没有准备的人分析数值算法，这当然不会是他首先想到的办法.第一个想法通常是前向误差分析：估计计算每一步引入的舍入误差，并设法累计这些误差在各步之间如何复合.

经验表明，对于数值线性代数中的大多数算法，前向误差分析比后向误差分析更难.事后看来，原因并不难解释.假设在计算机上使用一个久经检验的算法，例如求解 $Ax=b$.已知当 $A$ 病态时，得到的结果会始终不那么准确（见第 22 讲）.那么，前向误差分析如何捕捉这一现象呢？$A$ 的条件数是一个非常全局的性质，在求解 $Ax=b$ 时所涉及的单次浮点运算层面上几乎不可见.（下一讲会通过一个例子把这一点表现得更明显.）然而，前向分析无论如何都必须检测到这个条件数，才能最终得到正确结果.

简言之，对于大多数问题，最好的算法一般也不会做得比“对略有扰动的数据计算精确解”更好.后向误差分析正是一种与这种后向现实相适应的推理方法.

\section*{习题}

\begin{enumerate}[label=\textbf{15.\arabic*}]
    \item 下面每个问题都描述了一个在满足公理 (13.5) 和 (13.7) 的计算机上实现的算法.对每一个问题，判断算法是后向稳定、稳定但非后向稳定，还是不稳定，并证明结论，或者至少给出有说服力的理由.务必遵循正文中的定义.
    \begin{enumerate}[label=(\alph*)]
        \item 数据：$x\in\C$.解：$2x$，计算为 $x\oplus x$.
        \item 数据：$x\in\C$.解：$x^2$，计算为 $x\otimes x$.
        \item 数据：$x\in\C\setminus\{0\}$.解：$1$，计算为 $x\oslash x$.（满足式 (13.6) 的机器会给出完全正确的答案，但我们的定义基于较弱的条件 (13.7).）
        \item 数据：$x\in\C$.解：$0$，计算为 $x\ominus x$.（同样，真实机器可能比基于式 (13.7) 的定义表现得更好.）
        \item 无数据.解：$e$，从左到右使用 $\otimes$ 和 $\oplus$ 对
        \[
            \sum_{k=0}^{\infty}\frac1{k!}
        \]
        求和，直到某一项的大小小于 $\epsmach$ 时停止.
        \item 无数据.解：$e$，使用上述同一算法，但从右到左对级数求和.
        \item 无数据.解：$\pi$，穷举搜索区间 $[3,4]$ 中最小的浮点数 $x$，使得
        \[
            s(x)\otimes s(x')<0.
        \]
        这里 $s(x)$ 是在给定区间内稳定计算 $\sin(x)$ 的算法，$x'$ 表示浮点系统中紧随 $x$ 之后的下一个浮点数.
    \end{enumerate}

    \item 考虑计算矩阵完整 SVD 的算法.该问题的数据是矩阵 $A$，解是三个矩阵 $U$（酉矩阵）、$\Sigma$（对角矩阵）和 $V$（酉矩阵），满足
    \[
        A=U\Sigma V^*.
    \]
    这里说的是显式矩阵 $U$ 和 $V$，而不是由反射器乘积给出的隐式表示.
    \begin{enumerate}[label=(\alph*)]
        \item 解释这个算法后向稳定意味着什么.
        \item 事实上，由一个简单原因，这个算法不可能后向稳定.解释原因.
        \item 幸运的是，计算 SVD 的标准算法（第 31 讲）是稳定的.解释对于这类算法，稳定性意味着什么.
    \end{enumerate}
\end{enumerate}

\end{document}
